%%%%%%%%%%%%%%%%%%%%%%%%%%%%%%%%%%%%%%%%%%%%%%%%%%%%%%%%%%%%%%%%
\section{Reference Model and Evaluation Metrics}
\label{sec:case_study_reference_metrics}

The numerical references define the comparison basis for the scenarios in Section~\ref{sec:case_study_scenarios}.
The baseline and contact scenarios are compared with monolithic OpenModelica models of the corresponding closed-loop configuration.
The performance scenario uses a full-FEM co-simulation run as the high-fidelity benchmark.
All references are numerical comparison models.
No experimental ground truth is used.

The case-study scenarios are reproducible from the \syssimx{} repository.
The scenario notebooks define the co-simulation configurations used for the figures in this chapter.
The OpenModelica reference models belong to the controlled-pendulum Modelica package.
The reference trajectories are exported once and loaded by the plotting scripts.
This allows the thesis figures to be regenerated without recompiling the Modelica models.
Reproducing the exported reference trajectories requires OpenModelica \texttt{1.26.3}.

\paragraph{OpenModelica Reference Models.}
Table~\ref{tab:case_reference_models} lists the monolithic reference models used in the reported scenario evaluation.
Each reference contains the setpoint, controller, actuator, sensor blocks when required, and pendulum plant in one equation-based system.
Both references start at $0\,\mathrm{s}$ and use an output interval of $10^{-3}\,\mathrm{s}$.
The baseline reference stops at $2\,\mathrm{s}$.
The contact reference stops at $1\,\mathrm{s}$.

\begin{table}[tbp]
  \centering
  \small
  \caption{Monolithic OpenModelica reference models used in the case study.}
  \label{tab:case_reference_models}
  \begin{tabular}{p{4.0cm} p{5.0cm} p{5.2cm}}
    \toprule
    \textbf{Reference model} & \textbf{Scenario role} & \textbf{Main observables} \\
    \midrule
    \texttt{NoContact.Baseline}
      & Baseline closed-loop system without sensor quantization or contact
      & $\theta$, $\omega$, $u_{\mathrm{control}}$, $\tau$ \\
    \addlinespace
    \texttt{Contact.RigidContact}
      & Contact and multi-model switching scenario
      & $\theta$, $\theta_{\mathrm{meas}}$, \texttt{pid.I\_out}, contact flag \\
    \bottomrule
  \end{tabular}
\end{table}

\paragraph{Solver Configuration.}
The OpenModelica reference simulations use DASSL.
OpenModelica describes DASSL as an implicit BDF-based multistep solver with adaptive order and adaptive step-size control.
The exported reference trajectories use \(\texttt{startTime}=0\), \(\texttt{stopTime}=0.4\), \(\texttt{stepSize}=10^{-3}\,\mathrm{s}\), \(\texttt{tolerance}=10^{-6}\), and \(\texttt{solver}=\texttt{dassl}\).
The stored step size defines the output spacing of the reference trajectory.
DASSL selects its internal integration steps adaptively.

The Co-Simulation FMUs use CVODE as their default internal solver.
OpenModelica describes CVODE as a SUNDIALS solver for ordinary differential equation initial-value problems.
The exported stiff configuration uses variable-order and variable-step BDF integration with order 1 to 5.
The co-simulation macro step is set by the master algorithm.
It is separate from the internal FMU solver steps.

\paragraph{Full-FEM Benchmark Reference.}
The performance benchmark uses a full-FEM co-simulation run as its reference.
In this run, the FEM pendulum remains active over the full trajectory.
The switched \texttt{MasterPendulum} run is then compared against this high-fidelity benchmark.
The benchmark measures the computational benefit of switching relative to the trajectory deviation against full-FEM execution.

\paragraph{Compared Quantities.}
The main trajectory variable is the pendulum angle $\theta$.
For the baseline scenario, the angle is compared across multiple macro step sizes.
For the multi-model contact scenario, the contact event time and the PID integrator state are also evaluated.
For the performance scenario, wall-clock runtime and FEM-active time are reported in addition to the trajectory deviation.

\paragraph{Trajectory Error.}
All trajectory errors are evaluated on a common comparison grid.
The OpenModelica reference trajectory is interpolated to the accepted co-simulation output times.
For the pendulum angle, the pointwise error is
\begin{equation}
  e_{\theta,j}
  =
  \theta_{\mathrm{cs}}(t_j)
  -
  \theta_{\mathrm{ref}}(t_j) .
\end{equation}
The maximum angle error is
\begin{equation}
  E_{\infty,\theta}
  =
  \max_j |e_{\theta,j}| .
\end{equation}
The integrated angle error is
\begin{equation}
  E_{2,\theta}
  =
  \left(
    \frac{1}{T - T_0}
    \sum_j e_{\theta,j}^{2}\,\Delta t_j
  \right)^{1/2}.
\end{equation}
The maximum error identifies the worst trajectory deviation.
The integrated error captures the deviation over the full simulation interval.

\paragraph{Event-Time Error.}
For contact scenarios, the event-time error is evaluated separately from the trajectory error.
If $t_{\mathrm{event,cs}}$ is the event time detected by the co-simulation algorithm and $t_{\mathrm{event,ref}}$ is the corresponding event time in the OpenModelica reference, the absolute timing error is
\begin{equation}
  E_{t}
  =
  |t_{\mathrm{event,cs}} - t_{\mathrm{event,ref}}| .
\end{equation}
This metric is used only when a corresponding reference event can be identified unambiguously.

\paragraph{Switching Consistency.}
The multi-model contact scenario also checks whether runtime switching introduces discontinuities in the exposed plant state.
The \texttt{MasterPendulum} records every accepted switch in its \texttt{sync\_events} log.
For each switch, the deviations in \(\theta\) and \(\omega\) before and after state transfer are evaluated.
The maximum observed deviation is reported as the switching consistency error.
Values at numerical round-off level indicate that the model switch did not introduce a visible state jump.

\paragraph{Convergence Criterion.}
Numerical verification is based on macro-step refinement.
For the baseline scenario, the co-simulation is repeated with decreasing communication step sizes.
The trajectory is expected to approach the monolithic reference as \(\Delta t\) is reduced.
A scenario is interpreted as numerically verified when \(E_{\infty,\theta}\) and \(E_{2,\theta}\) decrease under refinement and the remaining deviations match the stated modeling differences.

\paragraph{Performance Criterion.}
The performance benchmark compares full-FEM execution with a switched configuration that activates FEM only near contact.
The full-FEM run defines the high-fidelity benchmark for this comparison.
The switched configuration is evaluated by its trajectory deviation against this benchmark and by its reduction in computational cost.
The computational benefit is reported as wall-clock speedup and reduction of FEM-active simulated time.